% !Mode:: "TeX:UTF-8"
\fancyhead[CO]{\song\wuhao ~\chaptername \quad 绪论~}
\fancypagestyle{plain}{\fancyhead[CO]{\song\wuhao ~\chaptername \quad 绪论~}}

\chapter{绪论}
% 第一节
\section{研究背景及意义}

Transformer 架构的大语言模型（Large Language Models, LLMs）近来在自然语言处理、计算机视觉、多模态学习等方向取得了重要进展。从 BERT、GPT 系列模型演变为现在的 LLaMA 和混合专家模型（Mixture of Experts, MoE），参数量和上下文长度在持续增加，这让深度学习任务的计算需求表现出爆炸式的增长。模型在推理部署过程中体现出的性能水平，会影响到实时交互体验及服务成本。研究证明，Transformer 模型的主要计算负载集中于矩阵乘（Matrix Multiplication, Matmul）、多头注意力机制（Multi-Head Attention，MHA）中。虽然矩阵乘在算力密集的硬件上已经有比较成熟的优化，但多头注意力由于涉及多个算子的链式组合并且具有高度复杂的访存模式，已经变成限制大模型推理延迟和吞吐量的主要瓶颈。

为加速注意力机制，行业出现了许多高效的底层优化方法，通过算子融合（Operator Fusion）和分块计算（Tiling），明显提高访存效率。不过，当前加速的实现主要依靠 NVIDIA 的 CUDA 生态系统和其专有硬件特点，包括张量核心（Tensor Core）、异步拷贝指令等。同时，国内外持续推出多种异构计算加速设备，海光 DCU 和寒武纪 NPU 等国产芯片的发展也很快。不同硬件架构在指令集、存储层次以及线程调度模型等方面有着较大区别，像 cuBLAS 和 rocBLAS 等现有厂商提供的底层高性能库，通常采用烟囱式的开发策略。因此开发者往往需要为特定平台花费大量时间进行代码编写和性能调优，这也带来了严重的软件生态碎片化问题，让算法跨平台迁移的成本和维护难度都变大了。

虽然解决这些问题的主要方式可以是自动化编译技术，但是在适配 Transformer 类算子和异构硬件时，现有的深度学习编译器仍然存在“语义鸿沟”和“控制力不足”的困难。例如 TVM 这类传统编译器大多使用基于循环的低级中间表示（Intermediate Representation, IR），它们所表示的张量计算的块级（Block）语义在降级（Lowering）的过程中丢失，这就使得编译器不容易分析得到线程块级别的优化机会。后来 Triton 这一类新型编译器虽然引入了块级抽象，但是它们主要依赖隐式的布局推断和黑盒化的数据搬运方式。面对海光 DCU 等拥有不同硬件特征和存储结构的国产硬件时，该隐式机制通常不能生成最优的访存指令，并且开发者难以通过显式方式干预数据通路来消除硬件瓶颈。所以，研究一种可以保留高层块级语义，同时显式解耦计算与数据搬运、以达到精细化硬件控制的算子生成新方法，成为现在需要解决的重要问题。

针对硬件异构化、软件生态碎片化、现有编译技术局限性的多重挑战，使得研究高效的跨平台算子自动代码生成技术具有较大的科学意义和应用价值。传统编译器在块级逻辑表达和底层访存控制上有所欠缺，本文通过多级中间表示（Multi-Level Intermediate Representation, MLIR）的多级抽象和方言（Dialect）机制，形成了一套可以把计算逻辑和数据搬运进行解耦的编译策略。通过引入显式的块级中间表示，在保留高层算子语义完整性的同时实现了对底层硬件特性的精确抽象和直接映射，并且使跨平台可移植性方面的异构硬件适配和扩展的工程复杂度得到明显降低。
本文通过自动代码生成技术替代高强度手工调优工作，解决了国产异构设备软件生态薄弱的问题。并设计通用编译优化 Pass，自动处理不同硬件架构下算子的访存冲突和流水线编排，从而快速生成特定硬件架构上的高性能算子。结合自动调优（Auto-tuning）技术，本文能够针对不同硬件参数构建参数搜索空间，利用剪枝策略自动搜索最优的执行配置。最终，在模型推理过程中，无论是在成熟的 NVIDIA GPU 还是新兴的国产加速卡上，本研究均能保持接近硬件峰值的算力利用率。


% \end{enumerate}
% 第二节
\section{国内外研究现状}

\subsection{Transformer 模型及其关键算子优化演进}

自从 Transformer\cite{vaswani2017attention} 架构被提出来之后，自然语言处理（Natural Language Processing, NLP）这个领域迎来了显著的发展变化。扩展定律\cite{kaplan2020scaling}（Scaling Law）表明，不断地堆叠 Transformer 层数、增加隐藏层维度并且扩大注意力头数，能够提高模型捕捉复杂语义的能力。为了能够得到更好的模型性能、更强的泛化能力，学术界和工业界一直遵循此定律。图 \ref{1_1} 展示了 Transformer 模型参数量随着时间变化的趋势，为了在逻辑推理、文本生成、代码补全等任务中取得质的突破，大模型从最初的 BERT-Base\cite{devlin2019bert} 模型发展到 GPT-3\cite{brown2020language}，再到后来万亿参数级别的混合专家模型\cite{cai2024survey}，模型参数量从亿级跨越到千亿级甚至万亿级，模型实现了指数级的增长。

\begin{figure}[htbp]
\centering
\includegraphics[width=1\textwidth]{1_1}
\caption{Transformer 模型参数量的增长趋势}\label{1_1}
\vspace{\baselineskip}
\end{figure}

然而随着模型规模持续上升，训练成本、推理成本急剧增长，这已然成为阻碍大模型大规模商业化应用的主要难题\cite{SCGD202404002}。面对这一挑战，近期新一代模型诸如 DeepSeek V3\cite{liu2024deepseek} 和 R1\cite{guo2025deepseek} 探索出低成本且高性能的新路径。这类模型并非单纯依赖盲目增加参数，而是借助高效的混合专家架构创新，并且还引入了强化学习\cite{yang2021exploration}（Reinforcement Learning, RL）驱动的后训练机制，进而使得模型的推理效率有了比较显著的提高。
这种策略的转变在算法方面降低了智能获取的边际成本，然而却对底层推理系统提出了更为严格的要求。对底层推理引擎进行极致优化，成为实现大模型技术普及的关键环节。
不论是复杂的稀疏激活策略，还是对超长上下文窗口的支持，均让在有限硬件资源条件下达成低延迟、高吞吐推理成为核心难题。


如图 \ref{1_2} 所示为模型推理时间占比随序列长度变化，在 Transformer 模型的推理过程中，计算负载主要集中在矩阵乘法与多头注意力机制这两个核心算子上。虽然矩阵乘占据了大量计算时间，但随着序列长度的增加，多头注意力机制\cite{vaswani2017attention}的计算复杂度与显存占用呈二次方增长趋势。多头注意力机制的计算也成为了制约长文本推理性能的重要瓶颈。
\begin{figure}[htbp]
\centering
\includegraphics[width=0.8\textwidth]{1_2}
\caption{Llama-2 模型推理时间占比随序列长度变化}\label{1_2}
\vspace{\baselineskip}
\end{figure}

传统的多头注意力机制的实现方式为“计算与访存分离”的方法。它首先计算查询向量（Query）与键向量（Key）的点积得到注意力分数矩阵，再将其写入高带宽全局显存中（HBM），然后再读取该矩阵进行 Softmax 归一化操作，最后再次写入显存。这种频繁的显存读写操作造成了严重的内存墙\cite{gholami2024ai}（Memory Wall）问题。GPU 的计算单元因长期等待数据而处于闲置状态，导致实际吞吐量远低于硬件理论峰值。

为解决这一问题，FlashAttention 系列算法\cite{dao2022flashattention, dao2023flashattention, shah2024flashattention, zadouri2026flashattention}被提出。它借助分块计算和重计算策略彻底改变了多头注意力机制的执行方式。首先，FlashAttention-v1\cite{dao2022flashattention} 把 GPU 片上高速共享内存（Shared Memory）作为缓存，运用算子融合技术在片上一次性完成矩阵乘、Mask 操作、Softmax、Dropout 等多个步骤。该方法防止中间大矩阵对全局显存反复读写，将访存复杂度从二次方降为线性。随后，FlashAttention-v2\cite{dao2023flashattention} 进一步优化了线程块方面的任务划分与并行策略，大量减少了非矩阵乘操作带来的同步开销。它通过优化线程束级的数据布局有效减轻了共享内存的 Bank 冲突。因此，多头注意力机制的计算在 NVIDIA Ampere 架构上的资源占用率有了明显的提高。其次，FlashAttention-v3\cite{shah2024flashattention} 针对 NVIDIA Hopper 架构\cite{nvidia_hopper}特性做了专门定制，它利用 WGMMA 矩阵指令集与张量存储加速器（Tensor Memory Accelerator, TMA）硬件特性\cite{luo2024benchmarking}，促使计算流水线与访存流水线的深度重叠\cite{soi2025optimal}，更接近硬件性能极限。最后，FlashAttention-v4\cite{zadouri2026flashattention} 又针对 NVIDIA Blackwell 架构\cite{nvidia_blackwell}设计了算法与算子（Kernel）流水线协同的注意力优化算法，它为了解决该架构上非矩阵乘单元成为性能瓶颈的问题，采用了更大分块（Tile）、软件模拟指数函数、条件性 Softmax 重缩放及 2-CTA 机制等技术，实现了性能进一步提升。

从 FlashAttention 系列的发展可以看出，大部分极致优化都较大程度上与 NVIDIA CUDA 指令集及其底层计算资源分配机制紧密相连。因为没有等效的高性能算子提供支持，如果将部署环境转换为海光 DCU 或者寒武纪 NPU 等国产异构平台，模型性能通常会出现急剧下降的情况。因此这种对于特定硬件生态的强烈依赖，致使 Transformer 模型在进行跨平台部署时，面临着高昂的代码重构成本、较大的维护压力。

虽然以 FlashAttention 作为代表性的手工深度优化算子，显著地提高 Transformer 模型的推理性能上限，然而其高昂的迁移成本、封闭的生态特性，已然无法契合当下异构计算硬件蓬勃发展的行业状况。当面临硬件架构快速更新换代、模型算子越来越复杂这两方面的压力时，仅仅依靠人工针对特定硬件去编写高性能算子的传统方式，已然难以为继\cite{cheng2025pipethreader, liu2026dr, liao2025kernelevolve}。究竟怎样能够摆脱对单一硬件厂商底层库的过度依赖，并且在确保算子具备极致性能的前提下，达成跨平台的自动化部署，这已然成为学术界与工业界共同探寻的关键要点。这也使得研究者们将视线投向了机器学习编译器技术\cite{JSJA202508003}，尝试借助自动化的代码生成、优化途径来打破这一困境。


\subsection{现有深度学习编译器的技术路径与局限}

深度学习模型的复杂度持续提高，若仅依靠手工方式来优化算子，在应对算子形态迅速演变、硬件后端碎片化\cite{lattner2021mlir}这两个问题时会面临较大困难。鉴于此，学术界和工业界着手开发了一系列深度学习编译器\cite{shi2023welder, ma2020rammer, rotem2018glow, wang2024ladder}，借助建立、变换中间表示的方式，达成从高层计算图到底层机器码的自动化生成。依据中间表示抽象层次的不同，当前的编译器技术主要分为两大流派：其中一种是以 XLA\cite{sabne2020xla}、TVM\cite{chen2018tvm} 作为代表的、基于标量循环的编译器，另一种是以 Triton\cite{tillet2019triton} 作为代表的、基于块级抽象的领域特定编译器。图 \ref{1_3} 展示了 TVM、Triton 和 DeepGen 在矩阵乘法算子上 IR 的对比情况，表 \ref{tab:table1} 展示了本文的工作与 TVM 及 Triton 类型编译器在核心设计方面的差异之处。虽然它们在特定阶段对技术发展起到了推动作用，然而当面对 Transformer 架构在异构平台上对于极致性能的需求时，均表现出不同程度的局限性。
\begin{figure}[htbp]
\centering
\includegraphics[width=1\textwidth]{1_3}
\caption{不同编译器框架下的矩阵乘法 IR 对比}\label{1_3}
\vspace{\baselineskip}
\end{figure}

\begin{table}[htbp]
\caption{TVM、Triton 与 DeepGen 的 IR 设计和优化能力对比}\label{tab:table1}
\vspace{0.5em}\centering\wuhao
\renewcommand{\arraystretch}{1.35}
\begin{tabular}{cccc}
\toprule[1.5pt]
比较维度 & TVM（TIR） & Triton（TTGIR） & DeepGen（DPIR）  \\\midrule[1pt]
设计目标 & 通用张量编译 & GPU 算子开发 & Transformer 算子生成 \\
IR 抽象粒度 & 循环级、标量级 & 块级张量 & 线程块级与线程级 \\
计算与访存关系 & 调度原语间接控制 & 高级原语统一描述 & Operation 显式解耦 \\
内存布局控制 & 依赖调度和张量化 & 依赖编译器推断 & LowerInfo 显式描述 \\
优化方式 & 循环变换与自动调度 & 块级 Pass 与启发式优化 & 规则 Pass 与自动调优 \\
跨平台特点 & 后端广，模板依赖强 & CUDA 生态成熟 & 统一 IR，Pass 可配置 \\
主要局限 & 块级语义易分散 & 细粒度控制较弱 & 通用算子覆盖待扩展 \\\bottomrule[1.5pt]
\end{tabular}
\vspace{\baselineskip}
\end{table}

第一代深度学习编译器（如 XLA、Halide\cite{ding2023hidet} 和 TVM 等）确立了计算与调度分离的核心设计理念。它们用多层嵌套循环描述计算逻辑，把高维张量运算降级为最基础的标量元素级读写操作。如图 \ref{1_3}(a) 所示，TVM 以张量表达式（Tensor Expression, TE）构建矩阵乘法 IR，TVM 的用户或自动调度器能够通过调度原语变换循环结构\cite{JSJA202503029}，像循环分块以适应 Cache 大小、循环重排优化数据局部性、向量化（Vectorize）和循环展开（Loop Unroll）等。AutoTVM\cite{chen2018learning} 和 Ansor\cite{zheng2020ansor} 引入机器学习技术，为了解决庞大调度空间里的参数选择难题，在 CPU 和早期 GPU 上实现了可与手工库媲美的性能。其方法主要是借助模板匹配（Pattern Matching）或草图生成（Sketch Generation），结合 XGBoost\cite{chen2016xgboost} 等算法自动搜索最优循环参数配置。

这种基于循环的表达方式通用性较高，却引发了比较严重的“语义丢失”情况。当编译器把矩阵乘法分解成 $for (i, j, k)$ 的标量循环后，原本矩阵计算的语义遭到了破坏。面对现代 GPU 架构时，硬件厂商为提高算力，引入了专用的矩阵加速单元。这些单元一般需要数据按照特定的分块形状和布局来加载。为了运用这些硬件特性，基于循环的编译器不得不执行复杂的张量化（Tensorization）过程，也就是通过模式匹配试着从破碎的标量循环中重塑矩阵语义，然后映射到硬件内建指令。这一过程较为脆弱，常常由于循环边界的细微变化或者复杂的算子融合逻辑而匹配失败，致使生成的代码退回到低效的标量计算策略，无法发挥硬件的峰值算力。

要解决标量循环难以自然映射到矩阵加速单元这一问题，新一代领域特定语言如以 OpenAI 推出的 Triton 为代表应运而生，很快成为 PyTorch\cite{paszke2019pytorch} 的默认后端技术。如图 \ref{1_3}(b) 所示，Triton 摒弃细粒度的标量循环描述，采用块级编程模型。在 Triton 里，基本操作单元不是单个浮点数，而是参数化的块。开发者用 tl.load、tl.dot、tl.store 等原语直接对张量块操作。这种设计让编译器能天然保留矩阵运算语义，更易映射到张量核心\cite{bhaskaracharya2020automatic} 等块级硬件指令上。同时，Triton 采用类似 Python 的语法，大幅降低高性能算子开发门槛，使算法工程师能以较高抽象层次编写 GPU 算子。

虽然 Triton 在 NVIDIA GPU 上有着出色表现，然而其设计理念里所包括的诸多隐式推断机制，在面临跨平台异构部署状况时，成为新的阻碍因素。在数据布局、访存\cite{zhou2026linear}这两方面，Triton 偏向于对共享内存的数据布局和存取模式进行自动管理。虽说它给出了 convert\_layout 这类接口，可具体的内存排布逻辑对于开发者而言仍然是半透明的。在 NVIDIA 架构环境下，这套启发式策略能够良好运行，不过当转移到海光 DCU 或者其他国产 AI 芯片上时，鉴于这些硬件有着不同的 Cache Line 大小、Bank 数量、冲突机制，Triton 的隐式策略常常难以生成最优的无冲突访问模式。
开发者很难借助位运算清晰地控制每一个数据的物理存储位置，致使性能在异构平台上出现比较严重的衰退情况。而在调度策略方面，Triton 的不少编译器优化 Pass 是依据 NVIDIA 硬件特性进行深度定制的。在移植到其他架构时，通常需要重新编写大量后端逻辑。

从 IR 设计理念上看，TVM 更强调通用性，其核心思想是将张量计算逐步降低为循环结构，并通过调度原语和自动调度系统在循环层面搜索高性能实现。该方式具有较强的后端适配能力，但在面向矩阵乘法和注意力机制等块级结构明显的算子时，块级计算语义容易在降级过程中被分散。后续若要映射到专用矩阵指令或复杂共享内存布局，通常需要依赖张量化匹配和较复杂的调度模板。

Triton 则更强调高性能 GPU 算子的开发效率，通过块级编程模型直接暴露张量块加载、矩阵乘和存储等原语，使开发者能够以较高抽象层次描述计算过程。该设计减少了手写 CUDA 算子的复杂度，并在 NVIDIA GPU 上取得了较好的工程效果。然而，Triton 对共享内存布局、线程映射和部分数据搬运路径采用较多隐式推断机制，开发者对底层访存模式的显式干预能力相对有限。当目标平台的线程组织、缓存层次或 Bank 冲突机制与 NVIDIA GPU 存在差异时，其优化效果更依赖后端实现和编译器内部启发式策略。

与上述两类路径相比，DeepGen 的设计目标不是替代通用编译器或通用 GPU 领域特定语言（Domain Specific Language, DSL），而是面向 Transformer 关键算子的跨平台高性能生成需求，在 IR 层显式区分线程块级计算语义和线程级执行细节。线程块级 IR 保留矩阵乘法、规约和逐元素计算等结构化语义，线程级 IR 则通过降级信息（LowerInfo）描述缓冲区（Buffer）在线程块、线程束和线程层级上的索引映射关系。该设计使 Software Pipeline、Local Split K、Block Swizzle 和 Access Scatter 等优化 Pass 能够直接作用于计算调度和访存路径，从而提高跨平台优化的可控性。相应地，DeepGen 当前也具有一定专用性，主要面向 Transformer 中计算密集和访存敏感的关键算子，对更广泛算子和动态控制流的支持仍需进一步扩展。

目前来看，现有编译方案在适配 Transformer 大模型、异构硬件的时候，会遇到三大核心挑战。首先是通用循环抽象难以有效利用张量核心，而高层块级领域特定语言又缺少对内存层次结构的显式控制。其次是现有框架过度依赖 CUDA 生态，结果在多样化的硬件后端很难实现统一的高性能。最后是算子融合优化导致调度搜索空间呈现组合爆炸，缺少领域知识引导的搜索很难收敛到高效解。基于这些挑战，本研究旨在建立一个具备表达力、可移植性、高效搜索能力的编译框架。

\subsection{硬件厂商高性能库的优势与局限}

在 PyTorch、TensorFlow\cite{abadi2016tensorflow} 等深度学习框架的底层实现环节里，硬件厂商所提供的高性能算子库起着基础性的关键作用。为了能在竞争激烈的硬件比拼当中通过基准测试，GPU 厂商、国产芯片厂商都投入了规模庞大的工程团队，依据特定的指令集架构、微架构特性，亲自手工编写并且进行优化核心数学库。这些库以 NVIDIA 的 cuBLAS\cite{cublas_doc}、cuDNN\cite{cudnn_doc}、 AMD 的 rocBLAS\cite{rocblas_doc}、MIOpen\cite{miopen_doc} 作为典型代表，展现出了当下通用矩阵乘与卷积（Convolution, Conv）运算的性能所能达到的最高水平。它们的高效之处在于对底层硬件进行了非常充分的利用：一方面，库开发者借助汇编级别的细致调整，直接采用底层汇编语言来编写算子\cite{tanner2018tensile}，从而能够精准地把控寄存器分配、指令流水线的排列、指令级并行。另一方面，库里面包括了成百上千个针对不同代际 GPU 架构、特定矩阵尺寸进行优化的算子变体，运行时系统会依据输入形状运用启发式算法动态挑选最优的实现方式。所以，在标准的矩阵乘法测试中，任何自动生成的代码要想超过这些经过数十年积累而成的手工库都是非常具有挑战性的。

\begin{figure}[htbp]
\centering
\includegraphics[width=0.80\textwidth]{1_4}
\caption{多头注意力机制算子的两种执行方式及推理时延}\label{1_4}
\vspace{\baselineskip}
\end{figure}

虽然厂商库在单个算子的性能方面展现出了出色的表现，然而在大模型快速迭代、异构部署的实际场景当中，其“黑盒”、“硬编码”的特性暴露出了较为严重的局限性。首先，不同硬件厂商的算子库存在着较为严重的生态隔离，进而形成了烟囱式的开发模式。这种隔离不只是体现在 NVIDIA 使用 CUDA 生态而海光 DCU 使用基于 ROCm 的 HIP 生态的底层差别上，还表现在 API 不兼容、函数签名、参数定义存在差异等问题上，致使将模型从 NVIDIA 平台迁移到国产硬件时需要重新编写大量胶水代码，产生了高昂的移植成本、对特定厂商生态的深度依赖\cite{DKJS202503002}。

再者，厂商库在算子融合能力方面存在着比较严重的瓶颈。图 \ref{1_4}(a) 展示了采用厂商提供的单算子实现和融合多头注意力机制的执行过程，其中绿色箭头代表从 HBM 中读取的操作，蓝色箭头代表写入的操作。在 Transformer 结构里，计算密集型算子和访存密集型算子交替出现，这就使得中间结果不得不反复写回到全局显存然后再重新读取，进而产生了巨大的带宽开销、“内存墙”问题。图 \ref{1_4}(b) 给出了未融合且反复访存 HBM 的多头注意力机制与融合型仅有一次 HBM 访存的多头注意力机制的时延比较。尽管部分厂商库开始提供有限的融合接口，然而这种支持是静态的且并非通用，难以灵活适应不断出现的新型激活函数\cite{fernandez2024telu}或者复杂的 Attention 变体\cite{lin2022survey}，致使编译器不得不退回到低效的逐一调用执行模式。

最后，在学术界提出新的低精度格式或者稀疏计算模式后，厂商库通常要数月甚至更长时间才能够提供官方支持。针对大模型里出现的变长序列、非标准形状等“长尾”场景，库函数的启发式选择常常失效，进而使得性能大幅下降，这是因为其优化资源常常集中在头部模型的标准尺寸方面。由于生态割裂、算子融合不足以及对模型演进响应滞后，即使硬件厂商的高性能库在标准算子上树立了性能标杆，也难以独自满足对于跨平台、端到端极致性能的追求。因此本研究致力于提出一种可以吸收手工库优化精华，同时拥有编译器灵活性的新型技术方案。

% 第三节
\section{论文主要研究内容}
针对 Transformer 类大模型在异构硬件部署中面临的适配复杂、手工优化成本高以及现有编译器表现不佳等挑战，本文设计了基于 MLIR 的 DeepGen 编译框架，实现高性能算子的跨平台自动生成，主要研究内容包括以下四个方面：

（1）面向异构硬件特性的块级分层中间表示设计与降级机制研究

为应对传统深度学习编译器于代码降级进程中出现的丢失算子块级语义这一状况，本文依托 MLIR 基础设施，设计出一套专用方言—— DeepGen。块级抽象建模方面，此 IR 策略运用分层设计理念，先在线程块方面针对矩阵计算进行粗粒度建模，将矩阵乘法与注意力机制的高维张量语义予以保留。接着，借助显式定义的 Alloc Buffer、Copy、Block、Matmul 操作（Operation），达成了计算逻辑与数据搬运的完全解耦。这种解耦设计让编译器能够单独对访存行为实施精细化管控，为消除异构硬件上的存储访问冲突创造了条件。自动化降级推断方面，针对从线程块级到线程级的映射难题，本文给出一种基于数据依赖关系的索引推断算法。该算法可按照算子的计算模式，自动推导线程块与线程束内的细粒度索引布局，达成从高层抽象至底层单指令多线程（SIMT）执行模型的自动化且语义无误的转换，有效屏蔽了底层硬件指令集的差别。

（2）异构平台下通用高性能编译优化 Pass 集合设计

鉴于 Transformer 关键算子在不同硬件上呈现出来性能差异的状况，本文构思了一套具备通用性且能够插拔的编译优化 Pass 集合，这套集合于计算调度、访存管理的各个方面均有所涉及。访存流水线优化方面，为对全局显存的高延迟加以掩盖，设计出了软件流水线（Software Pipelining）优化 Pass，借助多级缓冲技术达成数据的预取行为，以此促使计算指令与访存指令在时间维度上实现深度重叠。并行度增强优化方面，因为矩阵乘法与注意力机制里规约维度过大而致使线程块并行度不够这一问题，提出了局部规约轴拆分（Local Split K）优化 Pass。此项技术把规约维度的计算任务划分至多个逻辑线程组，并且在片上共享内存里对中间结果展开高效聚合，进而提高硬件占用率。访存冲突消除优化方面，面对不同硬件架构中存在的 Bank 冲突问题，提出了离散化访存（Access Scatter）与线程块重排（Block Swizzle）优化 Pass。通过针对访存索引施行离散化映射，明确地打散了连续地址的访问模式，在不改变计算语义的情形下消除了共享内存的 Bank 冲突，与此同时显著提高高速缓存（Cache）的命中率。

（3）基于空间剪枝策略的算子自动调优框架构建

针对手动调整编译参数时耗时费力且难以达到最优解的状况，本文搭建了一个闭环的自动性能调优框架。在参数空间生成方面，根据算子特性，抽象出固定种类的可进行调优的超参数，各算子根据自身计算特点构建该超参数配置文件，调优参数生成模块根据硬件特性组合生成基础调优空间，为剪枝模块提供数据空间基础。在高效剪枝策略方面，本文结合硬件特性设计了多重剪枝策略，例如静态分析去除会造成寄存器溢出或资源超限的无效配置，以及分析索引是否会造成 Bank 冲突等。最后框架通过实际运行并测量候选配置的性能，最终挑选出适配当前硬件特性的最优参数组合，达成了比肩厂商手工库性能的自动化代码生成。

（4）跨平台性能评估与端到端推理验证

为验证 DeepGen 框架有效性和泛化能力，本文于 NVIDIA A100/V100 GPU、海光 Z100/K100 DCU 这两类典型异构计算平台进行实验。
在算子级性能测试方面，本文将 DeepGen 与官方 BLAS 库以及当前主流编译器进行了多角度性能对比，并分析了不同实现的性能差异。
数值精度验证环节，本文采用计算结果 L2 范数相对误差，通过和标准 PyTorch 实现对比，以此证明本文提出的 IR 转换与优化策略在保持数值计算正确性上的可靠性。
端到端模型验证中，本文把生成的算子集成到 BERT、GPT-2、Llama-2 等主流 Transformer 模型推理流程里，验证优化后算子对模型整体推理延迟的提高效果。
此外，本文还对自动调优框架的开销进行实验分析，从候选配置数量、剪枝效果以及实际搜索时间等角度评估调优过程的效率，以说明该框架在降低手动调参成本的同时，能够在可接受的时间开销内搜索得到高性能执行配置。


% 第四节
\section{本文章节安排}
本文针对 Transformer 模型关键算子在异构平台上的自动化生成与性能优化展开研究，全文共分为五章，各章节的具体内容安排如下：

第 1 章绪论。本章首先讲述了以 Transformer 为代表的大语言模型在人工智能领域发展迅速，对底层算力需求极大，同时也揭示出当前异构硬件平台在算子开发时面临软件生态碎片化和适配成本高的问题。接着对国内外在深度学习编译器、高性能算子库方面的研究现状做了分析，尤其指出传统编译器存在块级语义表达缺失、异构兼容性差的情况，还有新兴领域特定语言工具存在隐式控制局限性。基于这些情况，确定了本文的研究动机、主要贡献、研究内容的整体框架。

第 2 章相关知识及理论介绍。本章为论文的研究建立起了必要的理论基础、背景知识。首先详细地介绍异构计算加速平台，着重对 NVIDIA GPU 和海光 DCU 的硬件架构特点及其编程模型展开对比分析。紧接着会深入剖析 Transformer 模型的网络架构，并且针对矩阵乘法与多头注意力机制的计算特性、访存瓶颈做量化分析。最后会介绍机器学习编译器的基本原理，着重说明 MLIR 框架的设计哲学、方言扩展机制、代码生成流程，以此论证其作为本文技术底座的合理性。

第 3 章 Transformer 模型关键算子的代码生成。此章为本文核心所在，细致说明了 DeepGen 编译框架的设计与实现详情。首先，给出一种面向硬件层次结构的多级 IR 设计，界定了块级与线程级的中间表示，还给出基于数据依赖关系的 IR 降级推断算法。其次，着重介绍针对异构硬件特性的通用优化 Pass 集合，详尽论述 Software Pipeline、Local Split K、Block Swizzle、Access Scatter 等关键优化技术的实现原理。最后，设计基于空间剪枝的自动调优框架，阐述参数空间的生成与剪枝策略，来达成最优编译配置的自动化搜索。

第 4 章实验结果。在这一章里，针对所提出的 DeepGen 框架展开全面的实验验证与性能评估工作。实验借助 NVIDIA GPU 和海光 DCU 这两种异构计算平台来进行，挑选不同规模的矩阵乘法与注意力机制算子当作测试对象，将其与厂商官方高性能库、主流编译器在多个维度上展开性能对比，进而分析 DeepGen 的性能优势。同时，本章还通过 Pass 消融实验评估各项优化策略对最终性能的影响，并从候选配置数量、剪枝效果和搜索时间等方面分析自动调优框架的成本与效率。此外，本章进行了数值精度验证工作，以及基于 BERT、Llama-2 等模型的端到端推理性能测试，有力地证实了本文方法具备有效性、跨平台可移植性。

第 5 章总结与展望。此章对全篇研究工作予以总结，归纳了本文于异构算子生成、中间表示设计、编译优化机制方面所获取的主要成果与创新之处。客观剖析了当下研究存在的局限，并且依据大模型技术的发展趋向，对未来研究方向作出了展望。
